\begin{abstract}
An expert demonstration is expensive. Its cost is human time, which compute does not reduce. A fixed budget therefore leaves one choice: what each
demonstration contains. Query-efficient interactive imitation learning decides
only when to call the expert. Two decisions are left over: which failure to
correct, and where the demonstration begins. We present \textbf{DISEIL}
(\textbf{D}emonstration d\textbf{I}stillation for \textbf{S}ample-\textbf{E}fficient
\textbf{I}mitation \textbf{L}earning). DISEIL keeps the same rule for calling the
expert and makes those two decisions. An episode runs to its end, and somewhere
in it the policy first becomes uncertain. An episode that does not achieve the
task is a failed episode. DISEIL groups a round's failed episodes into recurring modes.
For each mode it prescribes the starting configuration of a new
demonstration, checked for feasibility before any expert time is spent. We
evaluate across five tasks, with image and state observations, ten settings in
all, at a budget of 20 demonstrations, against the best baselines. DISEIL reaches
the highest mean success rate in every setting, on average about 3 percentage
points above the strongest baseline.
\end{abstract}
